\subsection{Recordatorios sobre los productos tensoriales de cuerpos.}
\label{subsection:IV.2.4.3-fr}

Para comodidad del lector, recordamos aquí algunas propiedades de los
productos tensoriales de cuerpos que utilizaremos en los números siguientes~;
para las demostraciones, remitimos a [1].

\begin{env}[4.3.1]
\label{IV.2.4.3.1-fr}
Recordemos que una extensión $L$ de un cuerpo $k$ se dice \emph{primaria} si la
mayor extensión algebraica separable de $k$ contenida en $L$ es el propio
$k$.
\end{env}

\begin{proposition}[4.3.2]
\label{IV.2.4.3.2-fr}
Sean $K$ y $L$ dos extensiones de un cuerpo $k$. Si $L$ es una extensión
primaria de $k$, entonces $\operatorname{Spec}(L\otimes_k K)$ es irreducible
y, si $\xi$ es su punto genérico, $\boldsymbol{k}(\xi)$ es una extensión primaria de $K$.
Recíprocamente, si para toda extensión finita separable $K$ de $k$,
$\operatorname{Spec}(L\otimes_k K)$ es irreducible, $L$ es una extensión
primaria de $k$.
\end{proposition}

\begin{proof}
Véase la demostración en [1], pp.~14-03 a 14-06.
\end{proof}

\begin{corollary}[4.3.3]
\label{IV.2.4.3.3-fr}
Si $k$ es separablemente cerrado \textup{(es decir,} si su clausura algebraica es
radicial sobre $k$\textup{)}, entonces, para dos extensiones cualesquiera $K$ y $L$ de
$k$, $\operatorname{Spec}(L\otimes_k K)$ es irreducible, y recíprocamente.
\end{corollary}

\begin{proof}
Esto se deduce inmediatamente de (4.3.2).
\end{proof}

\begin{corollary}[4.3.4]
\label{IV.2.4.3.4-fr}
Sean $L$ una extensión de un cuerpo $k$ y $L_s$ la mayor extensión algebraica
separable de $k$ contenida en $L$ \textup{(a veces llamada la \emph{clausura algebraica separable}
de $k$ en $L$)}~; supongamos que $L_s$ es de grado finito sobre $k$, y sea $K$
una extensión galoisiana de $k$ que contenga $L_s$. Entonces
$\operatorname{Spec}(L\otimes_k K)$ tiene $[L_s:k]$ componentes irreducibles,
cuya suma es el propio espectro, y los cuerpos residuales en los puntos
genéricos de esas componentes son extensiones primarias de $K$.
\end{corollary}

\begin{proof}
En efecto, se tiene $L\otimes_k K=L\otimes_{L_s}(L_s\otimes_k K)$ y se sabe que
$L_s\otimes_k K$ es isomorfo a un producto de $[L_s:k]$ cuerpos isomorfos a
$K$ (Bourbaki, \emph{Alg.}, cap.~VIII, \S~7, n\textsuperscript{o}~3,
cor.~2 del th.~1)~; por tanto, $L\otimes_k K$ es isomorfo a un producto de
$[L_s:k]$ anillos isomorfos a $L\otimes_{L_s}K$~; como $L$ es una extensión
primaria de $L_s$, la conclusión se sigue de (4.3.2).
\end{proof}

\begin{proposition}[4.3.5]
\label{IV.2.4.3.5-fr}
Sean $K$ y $L$ dos extensiones de un cuerpo $k$. Si $L$ es una extensión
separable de $k$, el anillo $L\otimes_k K$ es reducido. Recíprocamente, si para
toda extensión finita radicial $K$ de $k$, el anillo $L\otimes_k K$ es
reducido, $L$ es una extensión separable de $k$.
\end{proposition}

\begin{proof}
Para la demostración, véase Bourbaki, \emph{Alg.}, cap.~VIII, \S~7,
n\textsuperscript{o}~3, th.~1.
\end{proof}

\begin{corollary}[4.3.6]
\label{IV.2.4.3.6-fr}
Si $k$ es perfecto, entonces, para dos extensiones cualesquiera $K$ y $L$ de
$k$, $K\otimes_k L$ es reducido, y recíprocamente.
\end{corollary}

\begin{proof}
Esto se deduce inmediatamente de (4.3.5).
\end{proof}

\begin{corollary}[4.3.7]
\label{IV.2.4.3.7-fr}
Sea $L$ una extensión separable de $k$ y sea $K$ una extensión de $k$ tal que
$K$ o $L$ sea una extensión finita de $k$; entonces los cuerpos residuales del
anillo semi-local $L\otimes_k K$ son extensiones separables de $K$.
\end{corollary}

\begin{proof}
En efecto, se sabe que $L\otimes_k K$ es producto directo de esos cuerpos
residuales $E_i$~; para toda extensión $K'$ de $K$, $L\otimes_k K'$ es por
tanto isomorfo al producto directo de los anillos $E_i\otimes_K K'$~; como
$L\otimes_k K'$ es reducido, los $E_i$ son separables sobre $K$ en virtud de
(4.3.5).
\end{proof}

\begin{corollary}[4.3.8]
\label{IV.2.4.3.8-fr}
Si $K$ y $L$ son dos extensiones finitas separables de $k$, $K\otimes_k L$ es
producto de un número finito de extensiones separables de $k$.
\end{corollary}

\oldpage[IV-2]{59}

\begin{proposition}[4.3.9]
\label{IV.2.4.3.9-fr}
Si $k$ es algebraicamente cerrado, entonces, para dos extensiones cualesquiera
$K$ y $L$ de $k$, $L\otimes_k K$ es íntegro, y recíprocamente.
\end{proposition}

\begin{proof}
Es una consecuencia de (4.3.3) y (4.3.6), pues un cuerpo perfecto y
separablemente cerrado es algebraicamente cerrado.
\end{proof}

